% Focused framework presentation.
\section[Observable evolution]{Operators, duality, and statistical evolution}

\begin{frame}[t]{Evolution on the observable algebra}
\hypertarget{main:observables}{}
\begin{columns}[T,onlytextwidth]
  \begin{column}{.55\textwidth}
    \heading{Pull back the observable}
    On the observable algebra $\mathcal A=B_b(V)$, define
    \begin{equation*}
      \mathcal{K}_{s,t}f=f\circ\Phi_{s,t}.
    \end{equation*}
    The state map may be nonlinear, while its induced action on functions is linear.
    \begin{equation*}
      \begin{aligned}
        \mathcal{K}_{s,t}(fg)&=(\mathcal{K}_{s,t}f)(\mathcal{K}_{s,t}g),\\
        \mathcal{K}_{s,t}1&=1.
      \end{aligned}
    \end{equation*}
    \begin{equation*}
      \mathcal{K}_{s,r}\mathcal{K}_{r,t}=\mathcal{K}_{s,t}.
    \end{equation*}
  \end{column}
  \begin{column}{.40\textwidth}
    \vspace{7pt}
    \flowdiagram{pullback}
  \end{column}
\end{columns}
\slidecite{Manuscript \S\,1.1, transformation and evolution operators. State maps are Borel measurable.}
\end{frame}

\begin{frame}[t]{Observable evolution and probability evolution}
\hypertarget{main:duality}{}
For a bounded Borel observable $f$, the pushforward identity gives
\begin{equation*}
  \boxed{\langle f,\mu_t\rangle
  =\int f(\Phi_{s,t}(x))\mu_s(\dd x)
  =\langle\mathcal K_{s,t}f,\mu_s\rangle.}
\end{equation*}
\begin{columns}[T,onlytextwidth]
  \begin{column}{.48\textwidth}
    \heading{Schr\"odinger picture}
    Keep the observable fixed and evolve the law:
    \begin{equation*}
      \mu_t=(\Phi_{s,t})_\#\mu_s.
    \end{equation*}
    Integrate $f$ over the current distribution.
  \end{column}
  \begin{column}{.47\textwidth}
    \heading{Heisenberg picture}
    Keep the initial law fixed and pull back the observable:
    \begin{equation*}
      \mathcal K_{s,t}f=f\circ\Phi_{s,t}.
    \end{equation*}
    Integrate the evolved observable over $\mu_s$.
  \end{column}
\end{columns}
\medskip
Both descriptions give the same expectation. These identities also hold for
unbounded observables when the relevant integrals are finite.
\slidecite{Manuscript \S\,1.2.2, observable--measure duality and dynamical pictures.}
\end{frame}

\begin{frame}[t]{Infinitesimal evolution of observables}
\hypertarget{main:generator}{}
\heading{The deterministic generator on smooth test functions}
\begin{equation*}
  \mathcal L_t f=v_t\cdot\nabla f
  =\lim_{h\downarrow0}\frac{\mathcal K_{t,t+h}f-f}{h}.
\end{equation*}
The limit is understood on an appropriate smooth domain and in the chosen topology.
\heading{Differentiating the two time arguments}
For sufficiently regular $f$ and a smooth invertible flow,
\begin{equation*}
  \partial_t\mathcal K_{s,t}f=\mathcal K_{s,t}\mathcal L_t f,
  \qquad
  \partial_s\mathcal K_{s,t}f=-\mathcal L_s(\mathcal K_{s,t}f).
\end{equation*}
The terminal-time derivative acts on $f$ before pullback. The initial-time
derivative acts on the already pulled-back function.
\heading{Evolution of ensemble averages}
Under integrable domination along trajectories,
\begin{equation*}
  \frac{\dd}{\dd t}\langle f,\mu_t\rangle=\langle\mathcal L_t f,\mu_t\rangle.
\end{equation*}
\slidecite{Manuscript \S\,1.2.2, generators and observable evolution. \hyperlink{app:generator}{Time derivatives: Appendix A9.}}
\end{frame}

\begin{frame}[t]{Markov evolution extends the operator framework}
\hypertarget{main:markov}{}
\heading{Average over a transition law}
For a transition kernel $\mathsf K_{s,t}(x,\dd y)$, define
\begin{equation*}
  (\mathsf P_{s,t}f)(x)=\int f(y)\mathsf K_{s,t}(x,\dd y).
\end{equation*}
The same observable--measure pairing gives
\begin{equation*}
  \langle f,\mu_t\rangle=\langle\mathsf P_{s,t}f,\mu_s\rangle,\qquad
  \mathsf P_{s,r}\mathsf P_{r,t}=\mathsf P_{s,t}.
\end{equation*}
\heading{Which algebraic properties survive?}
Markov operators are linear, positive, and preserve the constant $1$.
They generally do not preserve products of observables.
\heading{Deterministic dynamics}
\begin{equation*}
  \mathsf K_{s,t}(x,\dd y)=\delta_{\Phi_{s,t}(x)}(\dd y)
  \quad\Longrightarrow\quad\mathsf P_{s,t}=\mathcal K_{s,t}.
\end{equation*}
In this special case, pullback also satisfies $\mathcal K(fg)=(\mathcal Kf)(\mathcal Kg)$.
\slidecite{Manuscript \S\,1.2.2, Markov kernels and deterministic evolution as a special case.}
\end{frame}

\begin{frame}[t]{Differential entropy under deterministic transport}
\hypertarget{main:entropy}{}
For a smooth positive density, define
\begin{equation*}
  S(\rho_t)=-\int\rho_t\log\rho_t\dd x,\qquad
  s_t^\rho=\nabla\log\rho_t.
\end{equation*}
Assume a classical continuity equation, finite integrals, justified differentiation,
and vanishing boundary terms.
\heading{Entropy rate}
\begin{equation*}
  \boxed{\frac{\dd}{\dd t}S(\rho_t)
  =\mathbb E_{\rho_t}[\nabla\cdot v_t]
  =-\mathbb E_{\rho_t}[v_t\cdot s_t^\rho].}
\end{equation*}
\begin{itemize}\setlength{\itemsep}{6pt}
  \item Positive mean divergence increases differential entropy.
  \item Negative mean divergence decreases it.
  \item Hamiltonian transport preserves entropy under these assumptions.
\end{itemize}
For the triangular example, $S(\rho_t)=S(\rho_0)+0.3t$ when the entropies are finite.
\slidecite{Manuscript \S\,1.2.4; Assessment 9.10. \hyperlink{app:entropy}{Entropy derivation: Appendix A13.}}
\end{frame}
